\chapter{异常安全、代码质量和工具链}

\section{问题入口：异常发生时对象还可信吗}

异常安全不是“是否使用异常”的争论，而是异常发生后程序状态是否仍然可信。看一个转账函数：

\begin{lstlisting}[language=C++,caption={中途失败会破坏状态}]
void transfer(Account& from, Account& to, int amount)
{
    from.withdraw(amount);
    to.deposit(amount);
}
\end{lstlisting}

如果 \code{withdraw} 成功，\code{deposit} 抛异常，钱从 \code{from} 扣了，却没有加到 \code{to}。这不是异常的问题，而是操作没有设计事务边界。异常安全要从不变量出发：函数失败后，哪些状态必须仍然成立？

\section{先把失败路径写成账本}

不要一上来背“基本保证、强保证”。先把转账函数的每一步列出来：

\begin{enumerate}
  \item 检查 \code{from} 余额是否足够。
  \item 从 \code{from} 扣钱。
  \item 给 \code{to} 加钱。
  \item 写审计日志或通知外部系统。
\end{enumerate}

如果第 \code{2} 步之后第 \code{3} 步失败，账户总额被破坏；如果第 \code{3} 步之后第 \code{4} 步失败，账户余额已经正确，但审计缺失。两个失败点的性质不同：前者破坏核心不变量，后者是外部副作用缺失。异常安全分析必须把“不变量”和“附加动作”分开。

一种更稳的转账接口，是让单个账户的余额修改本身不抛异常，所有可能失败的检查都在修改前完成：

\begin{lstlisting}[language=C++,caption={先验证再做不抛提交}]
class Account {
public:
    [[nodiscard]] bool can_withdraw(int amount) const noexcept;
    void withdraw_checked(int amount) noexcept;
    void deposit_checked(int amount) noexcept;
};

bool transfer(Account& from, Account& to, int amount)
{
    if (!from.can_withdraw(amount)) {
        return false;
    }

    from.withdraw_checked(amount);
    to.deposit_checked(amount);
    return true;
}
\end{lstlisting}

这个版本不是靠异常恢复，而是改变设计：提交阶段只做不会失败的内存内状态修改。真实系统里如果涉及数据库、远程服务、审计日志，就需要更强的事务机制；普通 \cpp 对象的异常安全不能自动覆盖外部世界。

\section{三种常见保证}

异常安全常用三个层级描述：

\begin{itemize}
  \item 基本保证：异常后对象仍然有效，不泄漏资源，但值可能改变。
  \item 强保证：异常后状态像函数没被调用过。
  \item 不抛保证：函数承诺不抛异常，常用 \code{noexcept} 表达。
\end{itemize}

例如 \code{std::vector::push_back} 在很多情况下提供强保证：如果扩容或元素构造失败，原向量仍保持原样。要做到这一点，通常需要先准备新资源，成功后再提交状态。

\section{先准备，后提交}

假设要更新配置对象，解析新配置可能失败。不要边解析边修改当前配置：

\begin{lstlisting}[language=C++,caption={边解析边修改会留下半更新状态}]
void reload_bad(Config& config, std::istream& input)
{
    input >> config.host;
    input >> config.port;
    input >> config.thread_count;
}
\end{lstlisting}

更好的做法是先解析到临时对象，全部成功后再赋值：

\begin{lstlisting}[language=C++,caption={先构造临时配置再提交}]
void reload(Config& config, std::istream& input)
{
    Config next;
    input >> next.host;
    input >> next.port;
    input >> next.thread_count;
    validate(next);
    config = std::move(next);
}
\end{lstlisting}

如果解析或校验失败，\code{config} 没有被修改。这就是强保证的基本形状。它通常会多一个临时对象，但换来更清楚的失败语义。

\section{强保证背后的普通代码模型}

强保证可以还原成三步普通代码：

\begin{enumerate}
  \item 在临时状态里完成所有可能失败的工作。
  \item 检查临时状态满足不变量。
  \item 用一个尽量不抛的操作提交新状态。
\end{enumerate}

配置热更新就是典型例子。坏版本直接改当前配置：

\begin{lstlisting}[language=C++,caption={半更新配置的失败路径}]
void reload_bad(Config& config, std::istream& input)
{
    input >> config.host;          // 可能已经改了 host
    input >> config.port;          // 这里失败
    input >> config.thread_count;  // 没执行
}
\end{lstlisting}

如果第二行失败，\code{config} 处在“新 host、旧 port、旧线程数”的混合状态。调用者不知道这个状态是否能继续使用。强保证版本把这种混合状态限制在局部变量里：

\begin{lstlisting}[language=C++,caption={临时状态承担失败风险}]
Config parse_config(std::istream& input)
{
    Config next;
    input >> next.host;
    input >> next.port;
    input >> next.thread_count;
    validate(next);
    return next;
}

void reload(Config& config, std::istream& input)
{
    Config next = parse_config(input);
    config = std::move(next);
}
\end{lstlisting}

现在失败只会发生在 \code{parse_config} 内部，原 \code{config} 还没被碰。如果 \code{Config} 的移动赋值可能抛异常，最后一步仍然有风险；这时可以把 \code{Config} 设计成用 \code{std::shared_ptr<const ConfigData>} 持有不可变数据，提交时只交换指针，或者提供不抛的 \code{swap}：

\begin{lstlisting}[language=C++,caption={用不抛 swap 提交}]
void reload_with_swap(Config& config, std::istream& input)
{
    Config next = parse_config(input);
    using std::swap;
    swap(config, next); // 应当设计为 noexcept
}
\end{lstlisting}

这不是模板口号，而是一条能检查的实现规则：失败风险集中在临时对象，公开对象只在最后一步切换。

\section{回滚对象适合局部状态，不适合所有副作用}

有时无法做到“提交前不修改”。可以用一个小的回滚对象保存旧值，失败时恢复：

\begin{lstlisting}[language=C++,caption={局部回滚对象}]
class BalanceRollback {
public:
    explicit BalanceRollback(Account& account)
        : account_{account}, old_balance_{account.balance()}
    {
    }

    void commit() noexcept
    {
        this->committed_ = true;
    }

    ~BalanceRollback()
    {
        if (!this->committed_) {
            this->account_.set_balance_for_rollback(this->old_balance_);
        }
    }

private:
    Account& account_;
    int old_balance_ = 0;
    bool committed_ = false;
};
\end{lstlisting}

这个模型只适合恢复内存内、可逆、恢复过程不抛异常的状态。已经发出去的网络请求、已经打印的日志、已经写入另一个进程的数据，不能靠析构函数“撤回”。所以异常安全不是魔法回滚，它要求你先判断哪些动作可逆，哪些动作必须放到提交点之后，哪些动作需要外部事务系统。

\section{析构函数不要抛异常}

析构函数常在栈展开期间执行。如果析构函数再抛异常，程序可能直接终止。因此析构函数应该吞掉或记录清理失败，而不是让异常逃出：

\begin{lstlisting}[language=C++,caption={析构中使用错误码清理}]
class TempFile {
public:
    ~TempFile()
    {
        std::error_code ignored;
        std::filesystem::remove(this->path_, ignored);
    }

private:
    std::filesystem::path path_;
};
\end{lstlisting}

如果清理失败必须报告，就提供显式 \code{close} 或 \code{commit} 函数，让调用者在可处理错误的位置调用。析构负责兜底清理，不负责复杂错误恢复。

\section{noexcept 是接口承诺}

\code{noexcept} 不是优化装饰，而是承诺。移动构造如果不抛，标准容器在扩容时更愿意移动元素：

\begin{lstlisting}[language=C++,caption={移动构造通常应 noexcept}]
Buffer(Buffer&& other) noexcept
    : data_{other.data_}, size_{other.size_}
{
    other.data_ = nullptr;
    other.size_ = 0;
}
\end{lstlisting}

不要随便给可能抛异常的函数标 \code{noexcept}。如果 \code{noexcept} 函数内部真的抛出异常，程序会调用 \code{std::terminate}。所以它必须和实现事实一致。

\section{编译警告和静态检查}

基础构建至少打开警告：

\begin{lstlisting}[language=CMake,caption={目标级警告选项}]
target_compile_options(my_target PRIVATE
    -Wall
    -Wextra
    -Wpedantic
    -Wconversion)
\end{lstlisting}

警告不是绝对真理，但它们是便宜的质量信号。比如未使用变量、隐式窄化转换、缺少返回值，经常能提前暴露真实 bug。不要把警告全局关掉，也不要为了“清零”写无意义代码。每个警告都要分类：真实问题、风格问题、工具误报，还是需要局部抑制。

\section{消毒器：让未定义行为早点爆}

地址消毒器可以发现越界、use-after-free 等问题：

\begin{lstlisting}[language=bash,caption={用消毒器构建测试}]
cmake -S . -B build-asan -DCMAKE_BUILD_TYPE=Debug \
  -DCMAKE_CXX_FLAGS="-fsanitize=address,undefined -fno-omit-frame-pointer"
cmake --build build-asan
ctest --test-dir build-asan --output-on-failure
\end{lstlisting}

线程消毒器可以发现数据竞争，但通常不能和地址消毒器一起使用。并发代码要单独跑：

\begin{lstlisting}[language=bash,caption={线程消毒器构建}]
cmake -S . -B build-tsan -DCMAKE_BUILD_TYPE=Debug \
  -DCMAKE_CXX_FLAGS="-fsanitize=thread -fno-omit-frame-pointer"
\end{lstlisting}

工具不是证明程序正确，而是扩大测试能抓到的问题范围。没有测试输入，消毒器也无事可做。

\section{代码审查清单}

自查一段 \cpp 代码时，可以按这个顺序看：

\begin{enumerate}
  \item 所有权：谁拥有资源，谁只是观察？
  \item 生命周期：引用、指针、视图是否可能悬空？
  \item 错误路径：异常或错误返回后对象不变量是否仍成立？
  \item 复制移动：是否有无意大对象复制？移动后对象是否被错误使用？
  \item 复杂度：主要循环是否重复扫描历史数据？
  \item 测试：空输入、单元素、边界值、失败路径是否覆盖？
\end{enumerate}

这个顺序不是固定流程，但它能避免只看语法风格。真正影响可维护性的，通常是所有权、错误边界和数据结构选择。

\begin{exercisebox}
选择前面 LRU 缓存项目，给每个公开函数写异常安全说明：容量为零、哈希表分配失败、Value 移动失败时，对象是否仍满足不变量？如果无法提供强保证，说明它至少能提供什么保证。
\end{exercisebox}
